\documentclass[a4paper,11pt]{article}

\newcommand{\TPnumero}{Project 6}
\newcommand{\macrosPath}{../../preamble/macros.tex}
% =========================================================
% Tutorial / lab preamble — Time Series Analysis module
% article class + fancyhdr + tcolorbox + listings (English)
% Author : Aymen Ben Brik (aymen.benbrik@esprit.tn)
% =========================================================

\usepackage[utf8]{inputenc}
\usepackage[T1]{fontenc}
\usepackage[english]{babel}
\usepackage{lmodern}
\usepackage{amsmath, amssymb, amsthm, mathtools, bm}
\usepackage{graphicx}
\usepackage{booktabs}
\usepackage{array}
\usepackage{multirow}
\usepackage{colortbl}
\usepackage{xcolor}
\usepackage{tikz}
\usetikzlibrary{positioning, arrows.meta, calc, shapes, fit, backgrounds, matrix}
\usepackage{listings}
\usepackage[most]{tcolorbox}
\usepackage{geometry}
\geometry{a4paper, margin=2cm, headheight=15pt}
\usepackage{fancyhdr}
\usepackage[hidelinks]{hyperref}
\usepackage{enumitem}

% --- Colours ---
\definecolor{espritBlue}{RGB}{30, 60, 110}
\definecolor{espritAccent}{RGB}{200, 80, 30}
\definecolor{boxEx}{RGB}{30, 60, 110}
\definecolor{boxQ}{RGB}{180, 120, 0}
\definecolor{boxC}{RGB}{30, 100, 60}
\definecolor{boxAtt}{RGB}{200, 80, 30}

% --- Listings ---
\definecolor{codebg}{RGB}{248,248,248}
\definecolor{codeframe}{RGB}{220,220,220}
\definecolor{keyword}{RGB}{0,82,155}
\definecolor{comment}{RGB}{88,110,117}
\definecolor{string}{RGB}{133,153,0}

\lstdefinestyle{pythonstyle}{
  language=Python,
  basicstyle=\ttfamily\small,
  keywordstyle=\color{keyword}\bfseries,
  commentstyle=\color{comment}\itshape,
  stringstyle=\color{string},
  backgroundcolor=\color{codebg},
  frame=single,
  rulecolor=\color{codeframe},
  frameround=tttt,
  breaklines=true,
  showstringspaces=false,
  tabsize=2
}
\lstdefinestyle{Rstyle}{
  language=R,
  basicstyle=\ttfamily\small,
  keywordstyle=\color{keyword}\bfseries,
  commentstyle=\color{comment}\itshape,
  stringstyle=\color{string},
  backgroundcolor=\color{codebg},
  frame=single,
  rulecolor=\color{codeframe},
  frameround=tttt,
  breaklines=true,
  showstringspaces=false,
  tabsize=2
}
\lstset{style=pythonstyle}

% --- Common macros (configurable path) ---
\providecommand{\macrosPath}{../../preamble/macros.tex}
\input{\macrosPath}

% --- Exercise counter ---
\newcounter{excounter}
\renewcommand{\theexcounter}{\arabic{excounter}}

\newtcolorbox[use counter=excounter]{exercise}[1][]{
  enhanced, breakable,
  colback=boxEx!5, colframe=boxEx, fonttitle=\bfseries,
  title={Exercise~\theexcounter\ifstrempty{#1}{}{ -- #1}},
  arc=2pt, boxrule=0.6pt
}
\newtcolorbox{question}[1][]{
  enhanced, breakable,
  colback=boxQ!5, colframe=boxQ, fonttitle=\bfseries,
  title={Question\ifstrempty{#1}{}{ -- #1}},
  arc=2pt, boxrule=0.5pt
}
\newtcolorbox{solution}[1][]{
  enhanced, breakable,
  colback=boxC!5, colframe=boxC, fonttitle=\bfseries,
  title={Solution\ifstrempty{#1}{}{ -- #1}},
  arc=2pt, boxrule=0.5pt
}
\newtcolorbox{warning}[1][]{
  enhanced, breakable,
  colback=boxAtt!5, colframe=boxAtt, fonttitle=\bfseries,
  title={Warning\ifstrempty{#1}{}{ -- #1}},
  arc=2pt, boxrule=0.5pt
}

% --- Headers / footers ---
\pagestyle{fancy}
\fancyhf{}
\fancyhead[L]{\textsc{\moduleNom} -- \auteurNom}
\fancyhead[R]{\TPnumero}
\fancyfoot[C]{\thepage}
\fancyfoot[L]{\auteurInstitution}
\fancyfoot[R]{\auteurEmail}
\renewcommand{\headrulewidth}{0.4pt}
\renewcommand{\footrulewidth}{0.2pt}

% --- Placeholder ---
\providecommand{\TPnumero}{Lab}


\title{Project 6: Mean--variance modelling and Value-at-Risk\\
\large ARMA + GARCH on a financial returns series}
\author{\auteurNom\\\auteurEmail\\\auteurInstitution}
\date{\anneeUniversitaire}

\begin{document}
\maketitle

\section*{Context}

This project closes the module by combining Chapters~3--6 in a single
volatility study. Starting from a real (or pre-supplied) financial
returns series, you will:
\begin{enumerate}
  \item fit a mean model (ARMA) on the level;
  \item test for ARCH effects on the residuals;
  \item fit a GARCH (or extended) model on the conditional variance;
  \item validate the joint specification through standardised
        residuals;
  \item produce a one-day Value-at-Risk and back-test it.
\end{enumerate}

The deliverable is a self-contained risk-management report.

\section*{Deliverables}
\begin{itemize}
  \item Reproducible \texttt{project6\_<name>.py} with seed $42$.
  \item A returns series of length $\geq 1500$ (daily). You may
        use the supplied \texttt{data/financial\_returns.csv} or
        bring your own (clear source citation).
  \item Report PDF (10--12 pages) including all figures, all
        numerical outputs, the back-test of the VaR.
  \item Optional 12-min oral defence.
\end{itemize}

% =========================================================
\begin{exercise}[Descriptive statistics and stylised facts]
\textbf{(1)} Plot the price series $P_t$ and the log-returns
$r_t = \log(P_t / P_{t-1})$.

\textbf{(2)} Report descriptive statistics on $r_t$: mean, standard
deviation, skewness, excess kurtosis. Compare with what would be
expected under a Gaussian model.

\textbf{(3)} Plot the sample ACF of $r_t$ and of $r_t^2$ up to
lag~$50$. Confirm the two stylised facts: returns are nearly
uncorrelated, but their squares are strongly autocorrelated.
\end{exercise}

% =========================================================
\begin{exercise}[Mean model: ARMA]
\textbf{(1)} On $r_t$, run ADF and KPSS to confirm stationarity.

\textbf{(2)} Identify candidate ARMA orders from ACF/PACF and fit
three candidates by MLE. Report AIC, BIC, Ljung--Box on residuals.

\textbf{(3)} Choose the best mean model. If no significant ARMA
structure is present, justify keeping a constant-mean model.
\end{exercise}

% =========================================================
\begin{exercise}[Detecting ARCH effects]
\textbf{(1)} On the residuals $\widehat\varepsilon_t$ of the chosen
mean model, run Engle's LM test at lags $5, 10, 20$. Report the
table.

\textbf{(2)} Plot the ACF of $\widehat\varepsilon_t^2$ and confirm
the visual evidence.

\textbf{(3)} Conclude: GARCH modelling is needed.
\end{exercise}

% =========================================================
\begin{exercise}[Variance model: GARCH$(1, 1)$]
\textbf{(1)} Fit the joint mean--variance model
$r_t = \mu_t + \varepsilon_t$, $\sigma_t^2 = \alpha_0 + \alpha_1
\varepsilon_{t-1}^2 + \beta_1 \sigma_{t-1}^2$ with Gaussian
innovations. Report all estimates with standard errors.

\textbf{(2)} Verify the stationarity condition $\widehat\alpha_1 +
\widehat\beta_1 < 1$ and report the implied unconditional
variance.

\textbf{(3)} Refit with a Student-$t$ innovation distribution and
compare AIC, BIC. Which one fits better? What does this say about
the tails of the data?
\end{exercise}

% =========================================================
\begin{exercise}[Diagnostics on standardised residuals]
\textbf{(1)} Compute $\widehat z_t = \widehat\varepsilon_t /
\widehat\sigma_t$ from the chosen GARCH.

\textbf{(2)} Plot the ACF of $\widehat z_t$ and $\widehat z_t^2$
up to lag $20$. Run Engle's LM test on $\widehat z_t^2$.

\textbf{(3)} Plot a QQ-plot of $\widehat z_t$ vs.\ the assumed
distribution. Comment on the marginal fit.
\end{exercise}

% =========================================================
\begin{exercise}[Volatility forecast and Value-at-Risk]
\textbf{(1)} On the chosen joint model, plot the in-sample
$\widehat\sigma_t$ together with the time series $r_t$. Identify
visually the most turbulent periods.

\textbf{(2)} Compute the multi-step variance forecast
$\widehat\sigma_{T+h}^2$ for $h = 1, \dots, 30$. Plot it together
with the unconditional variance.

\textbf{(3)} Compute the one-day Value-at-Risk at $\alpha \in
\{1\%, 5\%\}$, both conditionally
($\widehat\sigma_{T+1}$) and unconditionally
($\widehat\sigma$).
\end{exercise}

% =========================================================
\begin{exercise}[Bonus: VaR back-test]
On a hold-out test set (last $250$ business days), compute the
one-day rolling conditional VaR at $\alpha = 1\%$ and $\alpha =
5\%$. Count the number of \emph{exceedances} (days where the
realised loss exceeded the predicted VaR).
\begin{itemize}
  \item Expected number of exceedances at level $\alpha$ over $N$
        days: $\alpha N$.
  \item Run a Kupiec test (\texttt{statsmodels} or by hand) to
        check whether the realised exceedance count is consistent
        with the nominal $\alpha$.
\end{itemize}
\end{exercise}

% =========================================================
\section*{Marking grid (out of 100)}

\begin{center}
\begin{tabular}{p{8cm}rl}
\toprule
\textbf{Criterion} & \textbf{Points} & \\
\midrule
Descriptive statistics + stylised facts             & 10 \\
ARMA mean model + Box--Jenkins                      & 15 \\
Engle LM + diagnosis of ARCH effect                 & 10 \\
GARCH$(1, 1)$ Gaussian + Student-$t$ comparison     & 20 \\
Diagnostics on standardised residuals               & 15 \\
Forecast + Value-at-Risk                            & 15 \\
Bonus: VaR back-test (Kupiec)                       & 10 \\
Reproducibility + report quality                    &  5 \\
\midrule
\textbf{Total}                                      & \textbf{100} \\
\bottomrule
\end{tabular}
\end{center}

\bigskip
\hfill\auteurNom\par
\hfill\auteurEmail\par

\end{document}
